\documentclass[10pt]{article}
\usepackage[utf8]{inputenc}
\usepackage[T1]{fontenc}
\usepackage{amsmath}
\usepackage{amsfonts}
\usepackage{amssymb}
\usepackage[version=4]{mhchem}
\usepackage{stmaryrd}
\usepackage{graphicx}
\usepackage[export]{adjustbox}
\graphicspath{ {./images/} }

\begin{document}
нов $Q_{n}(z)$, то ряд Фурье по многочленам $\left\{K_{n}(z)\right\}$ функции $f(z)$, аналитической в односвязной области $G$ сходится к этой функции в среднем по площади области $G$ с весом $h(z)$, а при нек-рых дополнительных условиях и равномерно внутри области $G$.

Jum.: [1] S Z e g ö G., "Math. Z.", 1920, Bd 6, S. 167-202; 1921 , Bd 9, S. 167-90, 218-70; [2] C a r l e m a n T., "Ark. för mat., astr. och fys.", 1922-33, Bd 17, № 9, S. 1-30; [3] C е г е Г., ортогональные многочлены, пер. с англ., М., 1962; [4] Г е р он и м у с Я. Л., Многочлены, ортогональные на окружности и

\includegraphics[max width=\textwidth, center]{2023_07_08_2b7e4666f82d768d84e5g-1(1)}
нинградского физико-матем. об-ва», 1928, т. 2, в. 1, с. $155-79$; 185 ; [7] С уе т и н П. К., «Успехи матем. наук», 1966, т. 21, в. 2 , c. 41-88; [8] е г о же, «Тр. Матем. ин-та Ан СССР», 1971, т. 100 , с $1-96$.

ОРТОГОНАЛЬНЫЙ БАЗИС - система попарно ортогоналыных әлементов $e_{1}, e_{2}, \ldots ., e_{n}, \ldots$. гильбертова пространства $X$ такая, что любой элемент $x \in X$ однозначно представим в виде сходящегося по норме ряда

\[
x=\sum_{i} c_{i} e_{i}
\]

наз. рядом Ф у р ь е әлемен т а $x$ п о те м $\boldsymbol{e}\left\{e_{i}\right\}$. Обычно базис $\left\{e_{i}\right\}$ выбирается так, что $\left\|e_{i}\right\|=1$, и тогда он наз. о р т о н о р м и р о в а н ны м б а з и с о м. В этом случае числа $C_{i}$, наз. ко э фф ф

\includegraphics[max width=\textwidth, center]{2023_07_08_2b7e4666f82d768d84e5g-1(2)}
тон о рми ров ан ном у базису $\left\{e_{i}\right\}$, имеют вид $c_{i}=\left(x, e_{i}\right)$. Необходимым и достаточным условием того, чтобы ортонормированная система $\left\{e_{i}\right\}$ была базисом, является р а в н с т в С т е к ло о в а

\[
\sum_{i}\left|\left(x, e_{i}\right)\right|^{2}=\|x\|^{2}
\]

для любого $x \in X$. Гильбертово пространство, имеющее ортонормированный базис, является сепарабельным, и обратно, во всяком сепарабельном гильбертовом пространстве существует ортонормированный базис. Если задана произвольная спстема чисел $\left\{c_{i}\right\}$ такая, что $\Sigma_{i}\left|c_{i}\right|^{2}<\infty$, то в случае гильбертова пространства c базисом $\left\{e_{i}\right\}$ ряд $\Sigma_{i} c_{i} e_{i}$ сходится по норме к нек-рому элементу $x \in X$. Этим устанавливается изоморфизм любого сепарабельного гильбертова пространства пространству $l_{2}$ (т е о р е м Р

Лит.: [1] Ј ю с т е р н и к Л. А., С о б о ле в В. И., Элементы функционального анализа, 2 изд., М., 1965; [2] К о л м ого о о в А. Н., Фом и н С. В., Элементы теории функций и функционального анализа, 5 изд., М., 1981; [3] A х и е3 е р Н. И., Г л а з м а в И. М., Теория линейных операторов в гильбертовом пространстве, 2 изд., М., 1966. В. И. Соболев.

ОРТОГОНАЛЬНЫЙ ПРОЕКТӦР, о р т о п р о е кт о р,- отображение $P_{L}$ гильбертова пространства $H$ на его подпространство $L$ такое, что $x-P_{L} x$ ортогонально $P_{L} ! x-P_{L} x \perp P_{L} x$. 0 . п. есть ограниченный самосопряженный оператор, действующий в гильбертовом пространстве $H$, и такой, что $P_{L}^{2}=P_{L}$ и $\left\|P_{L}\right\|=1$. Обратно, если дан ограниченный самосогряженный оператор, действующий в гильбертовом пространстве $H$, и такой, что $P^{2}=P$, то $L_{P}=\{P x \mid x \in H\}$ является подпространством и $P$ есть $О$. п. на $L_{P}$. Два $О$. П. $P_{L_{1}}, P_{L_{2}}$ наз. о р т о г о на л ь н ы м и, если $P_{L_{1}} P_{L_{2}}=$ $=P_{L_{2}} P_{L_{1}}=0$; это эквивалентно условию, что $L_{1} \perp L_{2}$.

Свойства $О$. п.: 1) для того чтобы сумма $P_{L_{1}}+P_{L_{2}}$ двух О. п. была О. п., необходимо и достаточно, чтобы $P_{L_{1}} P_{L_{2}}=0$, в этом случае $\left.P_{L_{1}}+P_{L_{2}}=P_{L_{1}+L_{2}} ; 2\right)$ для того чтобы композиция $P_{L_{1}} P_{L_{2}}$ была О. п., необходимо и достаточно, чтобы $P_{L_{1}} P_{L_{2}}=P_{L_{2}} P_{L_{1}}, \quad$ в этом случае $P_{L_{1}} P_{L_{2}}=P_{L_{1} \cap L_{2}}$.

О. п. $P_{L^{\prime}}$ наз. ч а с т ь ю О. п. $P_{L}$, если $L^{\prime}$ есть подпространство $L$. IІри әтом $P_{L}-P_{L^{\prime}}$ является 0. п. на $L \div L^{\prime}-$ ортогональное дополнение к $L^{\prime}$ в $L$. В част-

\begin{center}
\includegraphics[max width=\textwidth]{2023_07_08_2b7e4666f82d768d84e5g-1}
\end{center}

Лит.: [1] Л ю с т е р ник Л. А., С о б о л е в В. И., Элементы функционального анализа, 2 изд., М., 1965; [2] A' х и езе у Н. И., Г л а з м а и. М., Теория линейных операторов в гильбертовом пространстве, 2 изд., М., 1966; [3] Р и с с Ф., С ёк е ф а л в в и- Н а д ь Б., Јекции по функциональному анализу, пер. с франц., 2 изд., М., $1979 . \quad$ В. И. Соболев.

ОРТОंГОНАЛЬНЫЙ РЯД - ряд вида

\[
\sum_{n=0}^{\infty} a_{n} \varphi_{n}(x), \quad x \in X
\]

где $\left\{\varphi_{n}(x)\right\}$ - ортонормированная система функций (онс) относительно меры $\mu(x)$ :

\[
\int_{X} \varphi_{i}(x) \varphi_{j}(x) d \mu(x)=\left\{\begin{array}{l}
0 \text { при } i \neq j, \\
1 \text { при } i=j
\end{array}\right.
\]

Начиная с 18 в. при изучении различных вопросов математики, астрономии, механики и физики (движение планет, колебание струн, мембран и др.) в исследованиях Л. Эйлера (L. Euler), Д. Бернулли (D. Bernoulli), А. Лежандра (A. Legendre), П. Лапласа (Р. Laplace), Ф. Бесселя (F. Bessel) и др. эпизодически появляются нек-рые специальные онс и разложения функций по ним. Определяюгцее же влияние на становление теории О. р. оказали:

a) исследования Ж. Фурье (Ј. Fourier, 1807-22) (Фурье метод решения краевых задач уравнений математич. физики) и в связи с ними работы Ж. Штурма I ЖЖ. Лиувилля (J. Sturm, J. Liouville, 1837-41); б) исследования П. Л. Чебышева по интерполированию и проблеме моментов (сер. 19 в.), повлекпиие за собой создание им общей теории ортогональных многочленов;

в) исследования Д. Гильберта (D. Hilbert, нач. 20 в.) по интегральным уравнениям, где, в частности, были установлены общие теоремы о разложөнии функций в ряд по онс;

г) создание А. Лебегом (Н. Lebesgue) теории меры и интеграла Лебега, придавшие теории О.р. современный вид.

Активному развитию теории 0 . р. в 20 в. способствует применение онс функций и рядов по ним в самых разнообразных разделах науки (математич. физика, вычислительная математика, функциональный анализ, квантовая механика, математич. статистика, операционное исчисление, автоматич. регулирование и управление, различные технич. задачи и т. п.).

Характерные результаты и направления исследований в теории 0. p.

\begin{enumerate}
  \item Пусть $X=\left[\begin{array}{ll}a, & b\end{array}\right], \quad d \mu(x)=d x$ - мера Лебега и $\left\{\varphi_{n}\right\}$ - онс. Тогда если $f \in L^{2}(a, b)$, то числа
\end{enumerate}

\[
a_{n}(f) \equiv\left(f, \varphi_{n}\right) \equiv \int_{a}^{b} f(x) \varphi_{n}(x) d x
\]

наз. к о әффи ци е т т м и Фу р ь е, а ряд (1) с $a_{n}=a_{n}(f)-р$ я д о м Ф у р в е функции $f$ по системе $\left\{\varphi_{n}\right\}$.

Система $\left\{\varphi_{n}\right\}$ з а м к н у т а относительно пространства $L^{2}$, если для любой функции $f \in L^{2}$ и любого числа $\varepsilon>0$ найдется полином

\[
\Phi(x)=\sum_{n=0}^{N} c_{n} \varphi_{n}(x)
\]

такой, что норма $\|f-\Phi\|_{2}<\varepsilon$. Система $\left\{\varphi_{n}\right\}$ п о л н а относительно $L^{2}$, если из условий $f \in L^{2}$ и $a_{n}(f)=0$ при всех $n \geqslant 0$ следует, что $f(x)=0$ почти всюду, т. е. $f-$ нулевоӥ элемент пространства $L^{2}$.

Если для нек-рой функции $f \in L^{2}$ выполняется равенство

\[
\sum_{n=0}^{\infty} a_{n}^{2}(f)=\int_{a}^{b} f^{2}(x) d x
\]

то говорят, что функция $f$ удовлетворяет у с л о в и ю за мк у тости Ля пунов а (Іли ра вен с т в П а р с в а л я). Это условие эквівалентно сходимости частных сумм ряда Фурье от $f$ по норме пространства $L^{2}$ к функции $f$.


\end{document}